\section{Uniform correction in complex quiver orbits}\label{sec:balancing}

Let $Q$ be a finite quiver, allowing loops and multiple arrows.
For Hermitian spaces $W_i$ of dimensions $m_i>0$, put
\[
 V=\bigoplus_{a:i\to j}\operatorname{Hom}(W_i,W_j),\qquad
 K=\prod_iU(W_i),\qquad G=\prod_iGL(W_i).
\]
We use the Hilbert--Schmidt metric on $V$ and the ordinary trace pairing
on the real vector space $\mathfrak a=\bigoplus_i\operatorname{Herm}(W_i)$.
The action and its Hermitian infinitesimal action are
\[
 (g x)_a=g_jx_ag_i^{-1},\qquad (C_xS)_a=S_jx_a-x_aS_i.
\]
Thus the vector field generated by $iS$ is $\xi_S(x)=iC_xS$.
Our real moment map and symplectic normalization are
\begin{equation}\label{eq:moment-normalization}
 \mu_i(x)=\sum_{h(a)=i}x_ax_a^*-\sum_{t(a)=i}x_a^*x_a,
 \qquad
 d\langle\mu,S\rangle=2\omega_0(\xi_S,\cdot),
 \quad \omega_0(iC_xS,C_xS)=\|C_xS\|^2.
\end{equation}
The common scalar acts trivially. A central covector is written
$\zeta=(\zeta_i I_{W_i})$ with $\sum_i m_i\zeta_i=0$.
A \emph{subrepresentation} is a collection $D_i\subset W_i$ with
$x_aD_{t(a)}\subset D_{h(a)}$ for every arrow.
A representation is \emph{module-simple} if its only
subrepresentations are zero and itself. This condition is independent
of a stability parameter and is stronger than having only scalar
endomorphisms.

Let $\Omega$ be a smooth closed $K$-invariant two-form on the Euclidean
ball $B_R\subset V$, and let $\rho:B_R\to\mathfrak a$ be an equivariant
moment map with
\begin{equation}\label{eq:perturbed-moment}
 \rho(0)=0,\qquad
 d\langle\rho,S\rangle=2\Omega(\xi_S,\cdot),\qquad
 \|\Omega-\omega_0\|\le\varepsilon.
\end{equation}
The last norm is the Euclidean operator norm. We do not assume that
$\Omega$ is of type $(1,1)$. Its indicated bound, with $\varepsilon<1$,
implies positivity on the pairs $(iC_xS,C_xS)$ used below.
We call $x$ $\rho$-balanced at $\zeta$ if $\rho(x)=\zeta$;
flat balancing means $\mu(x)=\zeta$.

\begin{theorem}\label{thm:uniform-balancing}
There are constants $A,\varepsilon_0>0$, depending only on
$M=\sum_i m_i$, with the following properties whenever
\eqref{eq:perturbed-moment} holds and $\varepsilon<\varepsilon_0$.
\begin{enumerate}
\item If $x\in B_{R/4}$ and $\mu(x)=\zeta$ is central, there is
$S\in\mathfrak a$ with $\|S\|_{\mathrm{op}}\le A\varepsilon$ and
$\rho(e^Sx)=\zeta$.
\item If $y\in B_{R/4}$ and $\rho(y)=\zeta$ is central, there is
$T\in\mathfrak a$ with $\|T\|_{\mathrm{op}}\le A\varepsilon$ and
$\mu(e^Ty)=\zeta$.
\item In a fixed complex orbit, the $\rho$-balanced representatives in
$B_R$ at the central value $\zeta$ form at most one $K$-orbit.
The corresponding local maps between balanced compact quotients are
continuous and inverse on nested neighborhoods. These conclusions
remain valid after restriction to a $G$-invariant analytic subset of $V$.
\end{enumerate}
No bound on a positive eigenvalue of $C_x^*C_x$ is required.
\end{theorem}

The proof extends the spectral-cut and degree argument of
\cite{Wue26a,Wue26c}. The projection estimate used in
\cite{Wue26d} is the appropriate error bound: both the coercive term and
the error contain the same crossing energy. Retaining the full stabilizer
in the argument also gives the reverse assertion and continuity when
stabilizers change.

\subsection{A relative estimate for every projection}
For a tuple $P=(P_i)$ of orthogonal projections, set
\[
 q_P(x)=\|C_xP\|^2
       =\sum_{a:i\to j}\|P_jx_a-x_aP_i\|^2.
\]
\begin{lemma}\label{lem:projection-error}
Under \eqref{eq:perturbed-moment},
\[
 |\langle\rho(x)-\mu(x),P\rangle|\le\varepsilon q_P(x)
 \qquad (x\in B_R).
\]
\end{lemma}
\begin{proof}
Write each arrow in its preserving and crossing blocks for $P$, and
write $x=n+u$ for the resulting orthogonal decomposition.
The circle $\exp(itP)$ fixes $n$. On its fixed vector space both
moment components paired with $P$ have zero differential, and their
normalizations at the origin are zero. Along $n+su$ we have
$\xi_P=s\,iC_uP$ and $\|C_uP\|=\|u\|$. Orthogonality keeps this segment
in $B_R$. The moment identities give
\[
 \langle\rho(x)-\mu(x),P\rangle
 =2\int_0^1s(\Omega-\omega_0)_{n+su}(iC_uP,u)\,ds.
\]
Its absolute value is at most $\varepsilon\|u\|^2=\varepsilon q_P(x)$.
\end{proof}

\subsection{The stabilizer and a strict spectral cut}
Fix $x$, let $H=K_x$, and define
\[
 \mathfrak h=\{S\in\mathfrak a:C_xS=0\},\qquad
 \mathfrak c=\{S\in\mathfrak a:Sh=hS\text{ for all }h\in H\},
 \qquad \mathfrak p=\mathfrak c\cap\mathfrak h^\perp.
\]
The commutation condition is understood vertexwise.
If $S\in\mathfrak p$, then $e^Sx\in V^H$. Any normalized equivariant
moment map takes its values on $V^H\cap B_R$ in $\mathfrak p$:
equivariance puts the value in $\mathfrak c$, and the vanishing
generating fields of $\mathfrak h$ show that its pairing with
$\mathfrak h$ is constant, hence zero, on this connected fixed ball.
In particular a central target attained at $x$ belongs to $\mathfrak p$.
Order the union of the spectra of all vertex matrices $S_i$,
including multiplicities. Adjacent eigenvalues below refer to this
single ordered list.

\begin{lemma}\label{lem:strict-cut}
If $S\in\mathfrak p$ and $\|S\|_{\mathrm{op}}=a>0$, there is a
global adjacent eigenvalue gap of size at least $a/(M-1)$.
For the vertexwise spectral projection $P$ above any strict global gap,
$q_P(x)>0$.
\end{lemma}
\begin{proof}
The identity tuple belongs to $\mathfrak h$, so $\operatorname{tr}S=0$.
The oscillation of its $M$ eigenvalues is at least $a$, proving the
gap estimate. A spectral projection $P$ commutes with $H$. If
$q_P(x)=0$, it also belongs to $\mathfrak h$. Since $I-P\in\mathfrak h$,
orthogonality gives
$\operatorname{tr}(SP)=\operatorname{tr}(S(I-P))=0$.
The mean eigenvalues above and below a strict gap cannot both vanish.
\end{proof}
For $M=1$ the action is trivial and all normalized moments vanish.
The following proof therefore assumes $M\ge2$.

\subsection{The forward and reverse degree arguments}
\begin{proof}[Proof of Theorem~\ref{thm:uniform-balancing}]
When $\varepsilon=0$, the moment identities and the common
normalization give $\rho=\mu$, so take zero correction. Assume
$\varepsilon>0$. We may take
\[
 A=2M,\qquad \varepsilon_0=\frac{\log2}{16A},\qquad a=A\varepsilon.
\]
If $\mathfrak p=0$, both moments at the starting point vanish and no
correction is needed. Otherwise consider, on the operator-norm ball of
radius $a$ in $\mathfrak p$,
\[
 F_s(S)=\mu(e^Sx)-\zeta+s(\rho-\mu)(e^Sx),\qquad 0\le s\le1.
\]
This homotopy takes values in the same fixed space $\mathfrak p$.
All its arguments have norm at most $e^{2a}\|x\|<R/2$.

At a boundary point diagonalize $S$ vertexwise and choose $P$ above
a gap $\delta\ge a/(M-1)$. An arrow entry with eigenvalue difference
$d$ has squared norm multiplied by $e^{2d}$. For an entry crossing
into $P$ the contribution to the paired moment increment is
$(e^{2d}-1)w$, with $d\ge\delta$; for one crossing out it is
$(1-e^{2d})w$, with $d\le-\delta$. Therefore
\begin{align}
 \langle\mu(e^Sx)-\mu(x),P\rangle
 &\ge(1-e^{-2\delta})q_P(x)
 \ge 2e^{-4a}\frac{a}{M-1}q_P(x), \label{eq:cut-lower}\\
 |\langle(\rho-\mu)(e^Sx),P\rangle|
 &\le\varepsilon e^{4a}q_P(x).                    \label{eq:cut-error}
\end{align}
Here $\delta\le2a$, and $P$ commutes with $S$.
Our constants make the first coefficient strictly larger than the
second. Lemma~\ref{lem:strict-cut} makes the crossing energy positive.
Thus $F_s$ has no boundary zero; no continuous choice of $P$ is needed.

The only zero of $F_0$ is $S=0$. Indeed
\[
 \frac{d}{dt}\langle\mu(e^{tS}x),S\rangle
   =2\|C_{e^{tS}x}S\|^2
\]
is strictly positive for $S\in\mathfrak p\setminus\{0\}$:
$C_{e^{tS}x}S=e^{tS}C_xS$, and $C_xS=0$ would imply
$S\in\mathfrak h\cap\mathfrak h^\perp$.
Moreover $DF_0(0)$ has positive quadratic form $2\|C_xS\|^2$ on
$\mathfrak p$, so its local degree is $+1$.
Brouwer degree gives a zero of $F_1$. This uses invertibility at each
fixed $x$, not a uniform lower eigenvalue estimate.

For the reverse assertion set $E=\rho-\mu$, define $\mathfrak p$
using $H=K_y$, and use
\[
 G_s(S)=\mu(e^Sy)-\mu(y)-sE(y).
\]
All terms again lie in $\mathfrak p$. Its initial zero is unique and
nondegenerate by the preceding calculation, which does not require
$\mu(y)$ to be central. On the boundary \eqref{eq:cut-lower} dominates
$|\langle E(y),P\rangle|\le\varepsilon q_P(y)$.
A zero at $s=1$ satisfies $\mu(e^Sy)=\rho(y)=\zeta$.
This proves existence of a flat balanced representative without an
independent polystability assumption.

For uniqueness suppose $y,gy\in B_R$ are $\rho$-balanced at the same
central value. Write $g=ke^T$. Equivariance and centrality reduce to
$y,e^Ty$. In an eigenbasis of $T$, the squared norm of $e^{sT}y$
is a sum of nonnegative exponentials, hence convex on $[0,1]$.
The entire path therefore stays in $B_R$, and
\[
 \frac{d}{ds}\langle\rho(e^{sT}y),T\rangle
 =2\Omega(iC_{e^{sT}y}T,C_{e^{sT}y}T)
 \ge2(1-\varepsilon)\|C_{e^{sT}y}T\|^2.
\]
Its endpoint values agree, so $C_yT=0$, and $gy=ky$.

Both corrections remain in any invariant analytic subset.
Uniqueness makes the compact-orbit maps independent of all choices.
For continuity choose convergent representatives and bounded correction
matrices. Every convergent subsequence of those matrices supplies a
balanced limit in the same complex orbit; uniqueness identifies its
compact orbit. The limiting matrix need not lie in the limiting
$\mathfrak p$. Compactness in the full fixed Hermitian algebra suffices.
The same proof gives parameter continuity for smooth families with a
common bound.

For explicit domains write
$\mathcal A_s=(Z\cap\mu^{-1}(\zeta)\cap B_s)/K$ and
$\mathcal B_s=(Z\cap\rho^{-1}(\zeta)\cap B_s)/K$.
The forward and reverse maps send radius $R/4$ into radius $R/2$
and are inverse wherever both are defined. In particular the forward
map identifies $\mathcal A_{R/8}$ with the open subset
$g^{-1}(\mathcal A_{R/8})$ of $\mathcal B_{R/4}$, where $g$ denotes
the reverse map; its image contains $\mathcal B_{R/16}$.
This proves the local statement with the required slack in radii.
\end{proof}

\begin{corollary}\label{cor:balancing-rate}
For a jointly smooth family satisfying
$\|\Omega_t-\omega_0\|_{B_{2r}}\le C(|t|+r^\alpha)$,
the correction on a fixed inner subball has norm
$O(|t|+r^\alpha)$, uniformly in stabilizer type and ratios of
nonzero arrow sizes. Smooth normalization gives $\alpha=1$.
The value $\alpha=2$ follows if the first spatial derivative vanishes,
for example under a verified inversion symmetry.
\end{corollary}
\begin{proof}
Apply Theorem~\ref{thm:uniform-balancing} on $B_{2r}$.
Taylor's theorem gives the two stated rates under their respective
derivative hypotheses.
\end{proof}
The theorem alone asserts neither differentiability across strata nor
a metric comparison. Both require identifying genuine tangent vectors
and removing complex gauge directions; we do this next.

